\documentclass[12pt,a4paper]{ctexbook}
\usepackage{geometry}
\geometry{margin=2.2cm}
\usepackage{setspace}
\setstretch{1.35}
\usepackage{hyperref}
\title{全宋诗选·poet.song.2000}
\author{古典文献整理}
\date{}
\begin{document}
\maketitle
\tableofcontents
\newpage
\section{緣識  其二九}
\textbf{宋太宗}\par\vspace{0.5em}
常思清凈世情閑，遥認天平我扣關。\par
莫道玄談些子是，鳳飛鶴宿九疑山。\par

\section{緣識  其三○}
\textbf{宋太宗}\par\vspace{0.5em}
乾坤運轉是尋常，人有依違自短長。\par
不覺貪生身外苦，如駒過隙百年光。\par

\section{緣識  其三一}
\textbf{宋太宗}\par\vspace{0.5em}
大丈夫兮兼文武，君子能仁潜喜怒。\par
强明非是盡周旋，滿眼真如不自取。\par
三春花卉正芬芳，發作癡迷我也狂。\par
道德修來長不斷，雲從龍變解舒張。\par
海浪無涯百川水，厚重高深難比擬。\par
今古猶來世界寬，忻聞爲善惡爲耻。\par
舟車有利便通津，斟量難事日其新。\par
勿笑我心孜孜化，貴爲人主不尊身。\par
共愛巔峰高岌岌，煩惱之中常急急。\par
先王典教是吾師，可以相宗後法則。\par
逍遥物外樂熙熙，松椿長養萬年基。\par
鵬翥引他鷦蝱怒，鳳皇爭與燕雀期。\par
心猨意馬須自縛，斯言之理還不若。\par
僻學邪聞似虛空，日月貞明閑銷鑠。\par
阿諛順旨察來情，瞽者忩忩不可行。\par
親親友愛慈敦睦，凡言俗態謾纏縈。\par

\section{緣識  其三二}
\textbf{宋太宗}\par\vspace{0.5em}
逍遥物外未曾拋，念誦時將玉磬敲。\par
雖不臨壇深發願，夢魂常憶大仙教。\par

\section{緣識  其三三}
\textbf{宋太宗}\par\vspace{0.5em}
居山城市罕逢人，傲世名高混世塵。\par
滿目煙霞常作伴，言知寒暑易爲春。\par

\section{緣識  其三四}
\textbf{宋太宗}\par\vspace{0.5em}
清潔恭虔禋上帝，郊儀大設康哉世。\par
人心肯悅化來同，禮樂咸修依古制。\par
景福照明何舉措，千官儼雅登雲路。\par
一陽初啟煦仁風，我運唯揚均徧布。\par
道德功慙情不已，常如取朽九重裏。\par
未曾一日貴爲尊，撫念蒼生皆赤子。\par
羽衛旌旛花作隊，雍熙民泰俱闤闠。\par
承平豈足盡言之，萬里山河觀掌內。\par

\section{緣識  其三五}
\textbf{宋太宗}\par\vspace{0.5em}
患脚法師不解走，少年心盡愛花柳。\par
爭知道味却無言，時得茶香全勝酒。\par

\section{緣識  其三六}
\textbf{宋太宗}\par\vspace{0.5em}
逍遥物外包相傳，遠望清虛碧似天。\par
莫道精修無報應，今生便見舊時緣。\par

\section{緣識  其三七}
\textbf{宋太宗}\par\vspace{0.5em}
凡碁妙手不可得，縱橫自在能消息。\par
不貪小利遠施張，舉措安詳求愛力。\par
曲須曲，直須直，打節斜飛防不測。\par
潛思靜慮一時間，取捨臨時方便逼。\par
牢己疆場煞三思，不驕不怯常翼翼。\par
勢輸他，勿動色，闇設機籌倍雅飾。\par
恒持自固最爲强，尤宜閑暇心先抑。\par

\section{緣識  其三八}
\textbf{宋太宗}\par\vspace{0.5em}
手談勝負觀其智，斂迹藏鋒不自制。\par
傲慢因從貪上虧，侵他理路無深計。\par
何齷齪，何凝滯，率爾圖謀太容易。\par
恃嶮揄揚似等閑，更翻劫子速求斃。\par
欺他不見心先喜，秪對前人失局勢。\par
已死更填隨手下，童蒙拙格一般藝。\par

\section{緣識  其三九}
\textbf{宋太宗}\par\vspace{0.5em}
十洲宮殿海漫漫，天上人間是可觀。\par
風裊長空清境界，凡夫興念大應難。\par

\section{緣識  其四○}
\textbf{宋太宗}\par\vspace{0.5em}
如來千百億，化身本是真。\par
駈駈煩惱不解嗔，我不能談說，舊却知新。\par

\section{緣識  其四一}
\textbf{宋太宗}\par\vspace{0.5em}
元正亨慶從此日，萬物咸新利貞吉。\par
送舊迎春一夜來，三元爲首衆情得。\par
萬彙潜臻俱是泰，玄化昭彰天作蓋。\par
氣和玉燭調陰陽，磨礱整頓知明代。\par
知明代，歌刑政，王道恢張千官盛。\par
任土陳儀四遠來，衣珮鏘鏘言不盡。\par
拱極之心文與武，乾坤澄濾煥今古。\par
丹鳳來儀見九重，我朝丕運還難遇。\par
感通玄貺時無比，五穀精華從茲始。\par
五色高吞日月明，百川盡歸朝宗水。\par

\section{緣識  其四二}
\textbf{宋太宗}\par\vspace{0.5em}
無名大道氣雄雄，日月昭明顯聖功。\par
舒卷好將浮世譽，亦同天上在人中。\par

\section{緣識  其四三}
\textbf{宋太宗}\par\vspace{0.5em}
逍遥聖境洞中天，五色雲藏大道仙。\par
七寶裝成銀世界，玉人皆是善心田。\par

\section{緣識  其四四}
\textbf{宋太宗}\par\vspace{0.5em}
重陽佳節菊花貴，婀娜如金葉青翠。\par
傍籬傍落一婆娑，發在秋天爲祥瑞。\par
引步登高甚奇見，想此儀形植霄漢。\par
恣情採折手中叉，滿袖冠簪香不散。\par
日飛黄鳥何老訝，看取叢叢光臺榭。\par
願細觀瞻向眼前，含珠露色應無價。\par
好是侵晨半開坼，風摇去住還難得。\par
清霜未布轉白雲，喜亦非常意南北。\par
柳岸煙沈和水淥，獨爾今朝看不足。\par
倏忽時間勝於春，隨分狂心無拘束。\par

\section{緣識  其四五}
\textbf{宋太宗}\par\vspace{0.5em}
心如明鑑照纖微，大道之中沒是非。\par
拂拭塵埃高挂閣，天邊日月有光輝。\par

\section{緣識  其四六}
\textbf{宋太宗}\par\vspace{0.5em}
道經念誦信無爲，清靜玄言好自持。\par
苦行精修心不倦，嚥津往往過齋時。\par

\section{緣識  其四七}
\textbf{宋太宗}\par\vspace{0.5em}
承平無事上元節，絲竹歌聲更互發。\par
自從雨後最晴明，不似往年今歲別。\par
死來無意在遨遊，常愧三光明皎潔。\par
千門萬戶樂喧喧，就中年少無蹔歇。\par
天街紅燄布羣星，晃煬人間闘明月。\par

\section{緣識  其四八}
\textbf{宋太宗}\par\vspace{0.5em}
上元時節皆相慕，皇都城裏萬家燈。\par
笙歌有處共誇稱，架肩疊足幾許層。\par
輕塵霧斂月明中，車騎雲軿五夜風。\par
意氣盡隨年少得，九衢填咽樂聲同。\par

\section{緣識  其四九}
\textbf{宋太宗}\par\vspace{0.5em}
木人莫把石牛騎，海變桑田故不知。\par
耕地種禾終是妄，真空相偶更無疑。\par

\section{緣識  其五○}
\textbf{宋太宗}\par\vspace{0.5em}
逢春酒一杯，百卉向陽開。\par
梨花看似梅，光揚迎日彩。\par
嫩柳引清風，黄鳥聲吹臺。\par

\section{緣識  其五一}
\textbf{宋太宗}\par\vspace{0.5em}
嶧陽之山，傳名曰桐，六律相沒五音中。\par
七軫弦調皆是意，審詳悞則亦如空。\par
清秋寂靜華堂深，聽之令我思沈吟。\par
幽蘭裏韻冬夜永，且合大道古人心。\par
舜製南風治蒼生，方知堯化廣聰明。\par
初彈將了移聲去，悲風秋思翻更互。\par
鳥啼別鶴何凄凉，秪聞斷續手揮忙。\par
慢引來催急風雨，寒泉妙滴散馨香。\par

\section{緣識  其五二}
\textbf{宋太宗}\par\vspace{0.5em}
識性低凡故不知，宣傳好事却爲非。\par
可憐此輩居人世，逆惡求真善不依。\par

\section{緣識  其五三}
\textbf{宋太宗}\par\vspace{0.5em}
蓬萊別有洞中天，摇曳祥雲五色鮮。\par
不是凡人常到處，精修要實在心專。\par

\section{緣識  其五四}
\textbf{宋太宗}\par\vspace{0.5em}
中秋八月祥風遍，慢上東輪雲遊轉。\par
清境將開照玉京，一時宮闕爲銀殿。\par
夜靜寒光半欹樓，空入簾帷白如練。\par
景致融怡露添寒，薄茶佳客隨情翫。\par
無問高下盡一般，爽氣周天如雪霰。\par
千家萬家仰天維，漏永更深方罷看。\par

\section{緣識  其五五}
\textbf{宋太宗}\par\vspace{0.5em}
黄芽藥就成金丹，理與參同契一般。\par
切是先求功行力，紅塵內見識應難。\par

\section{緣識  其五六}
\textbf{宋太宗}\par\vspace{0.5em}
玄元一氣不相監，清靜如存意馬銜。\par
大隠居鄽迷小道，絕巔頂上白雲巖。\par

\section{緣識  其五七}
\textbf{宋太宗}\par\vspace{0.5em}
牡丹花，深紅淺，忍把金刀枝上剪。\par
侵晨含露萬苞香，見者無非情展轉。\par
繁華影裏皆柔弱，密葉交加開閃爍。\par
朱欄別是一般春，蛺蝶悠颺閑自樂。\par
麗天和日景遲遲，思量往事不堪追。\par
風送園中來又去，豁然興感競移時。\par

\section{緣識  其五八}
\textbf{宋太宗}\par\vspace{0.5em}
妙手彈琴無向束，知之修鍊五音足。\par
先辨浮沈有指歸，弦頭制度相催促。\par
右手抑揚禁淫邪，左手徘徊堪瞻矚。\par
法於天，象於地，伏羲所造與心契。\par
先明理世見其真，六律含徽聲嘹唳。\par
從茲化被先賢慕，激濁揚清消喜怒。\par
太素仁風去住間，元和之氣皆徧布。\par
飛鳳在天不可測，大小龍吟不費力。\par
響應聽時有自然，舉措安詳能雅飾。\par
南風思政民俗化，順從平等無高下。\par
淳朴相傳今復興，逍遥道德後宗亞。\par
指要直掌須反善，拊安排齊似剪取。\par
聲來往，玄更玄，振兼文武情展轉。\par
古與今來千萬弄，幾人通達能妙用。\par
廣陵散好足仙蹤，胡笳十八堪鄭重。\par
堪鄭重，何清切，依憑伎倆能撥剌。\par
輕挑重打善間鉤，連蠲掄下輕微抹。\par
伯牙彈時如何美，汪汪洋洋似流水。\par
類例研究得剛柔，壞陵秋思無比擬。\par
敘志神和慢調⿰車爾，修身治性藏幽隠。\par
蒼龍鶴舞白雉飛，防奢止欲皆相準。\par

\section{緣識  其五九}
\textbf{宋太宗}\par\vspace{0.5em}
拊弄聲相引，兼能辨五音。\par
沈思調品切，句度更加吟。\par
緊慢弦中得，凝情指法深。\par
精英閑雅澹，停歇世途心。\par

\section{緣識  其六○}
\textbf{宋太宗}\par\vspace{0.5em}
全無體段弟一弱，十指乖張時更錯。\par
五音不辨信手彈，薄會人前言自樂。\par
弦不調，手不洗，更兼指下無道理。\par
吟爲失度故不知，把手空摇說宮徵。\par
合彈慢處却彈急，愚蒙堪笑情兀兀。\par
修持嚴潔甚難哉，岡知瀟灑心中出。\par
容易撫之何草草，顛來倒去直至老。\par
又耻人教却生嗔，虛度時光猶言好。\par
不聰人豈解，指法故然無。\par
拙見閑聲律，終朝自娛。\par
塵埃琴不飾，日月歲蕭疏。\par
慢共知音說，能彈好破除。\par

\section{迴文偈四首  其一}
\textbf{宋太宗}\par\vspace{0.5em}
上缺深意大誰教，指因緣善歸。\par

\section{迴文偈四首  其二}
\textbf{宋太宗}\par\vspace{0.5em}
會解垂要理，真詮羨非深。\par
意大誰教指，因緣善歸心。\par

\section{迴文偈四首  其三}
\textbf{宋太宗}\par\vspace{0.5em}
解垂要理真，詮羨非深意。\par
大誰教指因，緣善歸心會。\par

\section{迴文偈四首  其四}
\textbf{宋太宗}\par\vspace{0.5em}
重要理真詮。\par

\section{好字倒迴文}
\textbf{宋太宗}\par\vspace{0.5em}
好與煩性生，好聽。\par

\section{述字迴文偈七首  其一}
\textbf{宋太宗}\par\vspace{0.5em}
上缺留希佛背生。\par

\section{述字迴文偈七首  其二}
\textbf{宋太宗}\par\vspace{0.5em}
求比舊前憂，依述最亨知。\par
休以祐然留，希佛背生推。\par

\section{述字迴文偈七首  其三}
\textbf{宋太宗}\par\vspace{0.5em}
比舊前憂依，述最亨知休。\par
以祐然留希，紳背生推求。\par

\section{述字迴文偈七首  其四}
\textbf{宋太宗}\par\vspace{0.5em}
舊前憂依述，最亨知休以。\par
祐然留希佛，背生推求比。\par

\section{述字迴文偈七首  其五}
\textbf{宋太宗}\par\vspace{0.5em}
前憂依述最，亨知休以祐。\par
然留希佛背，生推求比舊。\par

\section{述字迴文偈七首  其六}
\textbf{宋太宗}\par\vspace{0.5em}
憂依述最亨，知休以祐然。\par
留希佛背生，推求比舊前。\par

\section{述字迴文偈七首  其七}
\textbf{宋太宗}\par\vspace{0.5em}
依述最亨知，休以祐然留。\par
希佛背生推，求比舊前憂。\par

\section{述字倒迴文三首  其一}
\textbf{宋太宗}\par\vspace{0.5em}
述依憂前舊，比求推生背。\par
以休知亨最。\par

\section{述字倒迴文三首  其二}
\textbf{宋太宗}\par\vspace{0.5em}
依憂前舊比，求推生背佛。\par
希留然祐以，休知亨最述。\par

\section{述字倒迴文三首  其三}
\textbf{宋太宗}\par\vspace{0.5em}
憂前舊比求，推生背佛希。\par
留然祐以休，知亨最述依。\par

\section{迴文偈二首  其一}
\textbf{宋太宗}\par\vspace{0.5em}
教訓大寬深，古羨協真齊。\par
要運解觀心，苦善業因迷。\par

\section{迴文偈二首  其二}
\textbf{宋太宗}\par\vspace{0.5em}

\section{好字拘就觀字正迴文}
\textbf{宋太宗}\par\vspace{0.5em}
道寬開訓新，生生業□□。\par

\section{太平興國七年季冬大雪賜學士}
\textbf{宋太宗}\par\vspace{0.5em}
輕輕相亞凝如酥，宮樹花裝萬萬株。\par
今賜酒卿時一盞，玉堂閑話道情無。\par

\section{試趙昌國}
\textbf{宋太宗}\par\vspace{0.5em}
秋風雪月天，花竹鶴雲烟。\par
詩酒春池雨，山僧道柳泉。\par

\section{賜陳摶  其一}
\textbf{宋太宗}\par\vspace{0.5em}
蒼生往世弊凋殘，今我如同赤子看。\par
大闡無爲三教盛，承平方說四夷寬。\par

\section{賜陳摶  其二}
\textbf{宋太宗}\par\vspace{0.5em}
餐霞成鶴骨，餌藥駐童顔。\par
靜想神仙事，忙中道路閑。\par

\section{賜陳摶  其三}
\textbf{宋太宗}\par\vspace{0.5em}
曾向前朝出白雲，後來消息杳無聞。\par
如今若肯隨徵召，總把三峰乞與君。\par

\section{賜蘇易簡  其一}
\textbf{宋太宗}\par\vspace{0.5em}
翰林承旨貴，清凈玉堂中。\par
應用咸依式，深嚴比更崇。\par
歸家思值日，入內集英風。\par
儒措門生盛，高明大化雄。\par

\section{賜蘇易簡  其二}
\textbf{宋太宗}\par\vspace{0.5em}
運偶昌時遠更深，果然穀在我中心。\par
從風臣偃光朝野，此日清華見翰林。\par
舉措樂時周禮法，思賢教古善規箴。\par
少年學士文明世，一守賢毫數萬尋。\par

\section{佛牙讚}
\textbf{宋太宗}\par\vspace{0.5em}
功成積劫印文端，不是南冊得恐難。\par
眼睹數重金色潤，手擎一片玉光寒。\par
鍊時百火精神透，藏處千年瑩采完。\par
定果熏修真祕密，正心莫作等閑看。\par

\section{送趙普}
\textbf{宋太宗}\par\vspace{0.5em}
忠勤王室展宏謨，政事朝堂頼秉扶。\par
解職暫酬卿所志，休教一念遠皇都。\par

\section{句  其一}
\textbf{宋太宗}\par\vspace{0.5em}
欲餌金鈎深未到，磻溪須問釣魚人。\par

\section{句  其二}
\textbf{宋太宗}\par\vspace{0.5em}
薄德終慚舉，通才例上居。\par

\section{句  其三}
\textbf{宋太宗}\par\vspace{0.5em}
時向玉堂尋舊迹，八花磚上日空長。\par

\section{句  其四}
\textbf{宋太宗}\par\vspace{0.5em}
鑾輿臨紫塞，朔野凍雲飛。\par

\section{句  其五}
\textbf{宋太宗}\par\vspace{0.5em}
好事盡輸純與直，謾勞頰舌湧如泉。\par

\section{句  其六}
\textbf{宋太宗}\par\vspace{0.5em}
奉使南遊多好景。\par

\section{句  其七}
\textbf{宋太宗}\par\vspace{0.5em}
君臣千載遇。\par

\section{句  其八}
\textbf{宋太宗}\par\vspace{0.5em}
二三千客裏成事，七十四人中少年。\par

\section{句  其九}
\textbf{宋太宗}\par\vspace{0.5em}
珍重老臣純不已，我慚寡昧繼三皇。\par

\section{句  其一○}
\textbf{宋太宗}\par\vspace{0.5em}
一箭未施戎馬遁，六軍空恨陳雲高。\par
鑾輿曠紫塞，朔野陣雲飛。\par

\section{上方}
\textbf{陳省華}\par\vspace{0.5em}
四望平川獨一峰，峰前瀟灑是蓮宮。\par
松聲竹韻千年冷，水色山光萬古同。\par
客到每憐樓閣異，僧言因得鬼神功。\par
懸民遥喜行春至，鼓腹閑歌夕照中。\par

\section{非風幡動仁者心動頌}
\textbf{釋本先}\par\vspace{0.5em}
非風幡動唯心動，自古相傳直至今。\par
今後水雲人欲曉，祖師直是好知音。\par

\section{見色便見心頌}
\textbf{釋本先}\par\vspace{0.5em}
若是見色便見心，人來問著方難答。\par
更求道理說多般，孤負平生三事衲。\par

\section{明自己頌}
\textbf{釋本先}\par\vspace{0.5em}
曠大劫來祇如是，如是同天亦同地。\par
同地同天作麽形，作麽形兮無不是。\par

\section{薔薇詩一首十八韻呈東海侍郎徐鉉}
\textbf{李從善}\par\vspace{0.5em}
綠影覆幽池，芳菲四月時。\par
管弦朝夕興，組繡百千枝。\par
盛引牆看徧，高煩架屢移。\par
露輕濡綵筆，蜂誤拂吟髭。\par
日照玲瓏幔，風摇翡翠帷。\par
早紅飄蘚地，狂蔓挂蛛絲。\par
嫩刺牽衣細，新條窣草垂。\par
晚香難暫捨，嬌態自相窺。\par
深淺分前後，榮華互盛衰。\par
尊前留客久，月下欲歸遲。\par
何處繁臨砌，誰家密映籬。\par
絳羅房燦爛，碧玉葉參差。\par
分得殷勤種，開來遠近知。\par
晶熒歌袖袂，柔弱舞腰支。\par
膏麝誰將比，庭萱自合嗤。\par
勻妝低水鑑，泣淚滴烟𩅰。\par
畫擬憑梁廣，名宜亞楚姬。\par
寄君十八韻，思拙媿新奇。\par

\section{寄贈終南英公上人}
\textbf{臧丙}\par\vspace{0.5em}
是箇人言好性靈，鎬京碑記念千廳。\par
五言出格爲詩匠，百盞長杯應酒星。\par
曾把篆求身上紫，幾將金買面前青。\par
多聞國士相尋訪，莫把松門畫夜扃。\par

\section{閒居}
\textbf{林緒}\par\vspace{0.5em}
宿草江邊閣，臨流發嘯歌。\par
所觀得病少，豈在蓄書多。\par
生子能無貴，居身不可過。\par
苟圖銷旦夕，青鬢長如何。\par

\section{送僧歸天甯萬年禪院}
\textbf{安德裕}\par\vspace{0.5em}
跡自青門遠，田衣賁在躬。\par
舊房千嶠外，歸棹五湖東。\par
地力薑畦沃，年支芋盎充。\par
從支乃榮道，一與祖心同。\par

\section{詩一首}
\textbf{馬文斌}\par\vspace{0.5em}
希奇寶象獸中王，猛毅雄心世不當。\par
四足端然如玉柱，雙牙利劍若金鋩。\par
立觀峭峻成山岳，動必摇形見者慌。\par
但以聲名告醜類，從今何敢作災殃。\par

\section{花雨比下秦中}
\textbf{田錫}\par\vspace{0.5em}
雨裏飛花片片紅，雨微花亂轉溟濛。\par
川原何處連天草，簾幕無人半日風。\par
寂寞劉楨新病後，淒迷莊舄苦吟中。\par
可堪更是黄昏景，燕冷鶑寒恨想同。\par

\section{寄蒲城宋白小著}
\textbf{田錫}\par\vspace{0.5em}
春日閑銷一局棋，春愁還得數篇詩。\par
高吟大醉何人問，英略雄圖久已知。\par
破産雖無容足所，丈夫豈合以家爲。\par
時來富貴終須有，懶學梁鴻賦五噫。\par

\end{document}